\section{YuE2: prediction-only transfer to locked recordings}
\label{sec:yue2locked}

The extension generated 500 recordings from the shared frozen Muse prompts.
Of these, 498 supported a native 30-second view: 398 development recordings
and 100 locked recordings. Two shorter recordings were excluded from this
Native30 analysis, not padded into eligibility. The experiment reported here
uses the 100 locked AI recordings only and does not use the expanded models.

\subsection{Fixed model catalogue and protocol deviation}
All 255 nonempty subsets of the eight pre-existing descriptor families were
evaluated with five previously fitted ordinary-group-holdout models per subset.
These models use the original cohort, the all-quantity setting, values plus
missingness, and a fixed decision threshold of 0.5. Training medians, scaling
and coefficients were reused, with zero new fits. The preflight checked
disjoint item, group and component identities between training and the locked
recordings. Each prediction was replayed deterministically. The result has
1,275 models and 127,500 prediction occurrences, but only 100 test recordings.

The originally proposed order was to score the held-out generator before
expanded development. Actual locked scoring occurred afterwards. This is a
protocol-order deviation, and the present report does not retrospectively
claim prospective preregistration. The exhaustive scoring catalogue was fixed
before this scoring run; no subset was chosen as a winner from these outcomes.

\subsection{Sensitivity, not two-class accuracy}
For a fixed model, sensitivity is the fraction of these 100 AI recordings
predicted as AI. There are no human test recordings in this experiment, so
specificity, balanced accuracy and ROC AUC are not estimable. A trivial
always-AI predictor would also achieve perfect sensitivity.
Table~\ref{tab:yue2locked} reports the single families and two pre-existing
combinations; the exhaustive 255-subset table is retained separately.

\begin{table}[htbp]
\centering
\caption{Locked YuE2 sensitivity (percent). Means and ranges are descriptive
across five old models evaluated on the same recordings, not confidence
intervals or independent test-set repetitions.}
\label{tab:yue2locked}
\begin{tabular}{lrr}
\toprule
Family/subset & Mean & Range \\
\midrule
S & 94.8 & 93--97 \\
D & 62.6 & 62--65 \\
R & 84.8 & 84--85 \\
P & 28.0 & 28--28 \\
F & 87.8 & 85--90 \\
H & 87.0 & 85--90 \\
SC & 90.2 & 88--91 \\
BC & 66.8 & 29--100 \\
S+D+R+P & 96.0 & 95--97 \\
S+D+R+P+F+H+SC+BC & 96.2 & 95--97 \\
\bottomrule
\end{tabular}
\end{table}

The four-family and eight-family combinations have similar descriptive
sensitivity, but this does not establish equal overall performance or an
incremental benefit from the additional families. BC alone varies strongly
across the old models. It must be assessed jointly with human false-positive
rates before interpreting utility. The results do not identify the acoustic
cause of the label association and cannot establish an intrinsic AI signature.

\subsection{Reproduction and retained evidence}
The code repository is
\url{https://github.com/ezmonyi/ai-music-detection}.
The versioned scripts are
\path{current/yue2_500_extension_20260911/preflight_locked_yue2_v1.py},
\path{score_locked_yue2_v1.py}, and \path{report_locked_yue2_v1.py} in that
same directory. The scoring protocol and freeze are retained alongside them.
The source scoring COMMIT SHA256 is
\path{b43744c0775c14d040579537316b40d1c6a3479649afc285faa176e76abf87f5}.
All 1,276 committed products were verified after local recovery.

Under \path{/Users/yi/Documents/report/music_ai_detection_20260912/},
\path{current_results/locked_yue2_scoring_v1/} contains individual predictions
and \path{current_results/locked_yue2_report_v1/all_255_combinations.csv}
contains the complete descriptive comparison. These paths identify preserved
results, not a claim that all remaining experiment arms are finished.
